%\ifodd\value{page}\hbox{}\newpage\fi
\chapter*{Śrī Lalitā Sahasranāma — Śloka 8}
\addcontentsline{toc}{chapter}{Śloka 8 - Nomi 18,19}

\begin{sloka}
{\sanskritfont
\begin{tabular}{c}
कदम्बमञ्जरीकॢप्तकर्णपूरमनोहरा ।\\
ताटङ्कयुगलीभूततपनोडुपमण्डला ॥ ८॥
\end{tabular}

\vspace{1.5em}

{\middlesize \iastfont
kadamba-mañjarī-kḷpta-karṇapūra-manoharā |\\
tāṭaṅka-yugalī-bhūta-tapanōḍupa-maṇḍalā ||8||}

\vspace{1em}

{\middlesize
\anvaya{%
कदम्बमञ्जरीकॢप्तकर्णपूरमनोहरा ।
ताटङ्कयुगलीभूततपनोडुपमण्डला ।
}
}

\middlesize \iastfont \itmean{%
«Incantevole per l’ornamento auricolare formato da un grappolo di fiori di \textit{kadamba};  
i cui due orecchini sono divenuti i dischi del sole e della luna.»%
}

}
\end{sloka}

\wordmeaning{%
\wordline{कदम्ब-मञ्जरी-कॢप्त-कर्णपूर-मनोहरा}{kadamba-mañjarī-kḷpta-karṇapūra-\\ manoharā}
  { incantevole, con un \textit{karṇapūra}, ornamento che “riempie” l’orecchio, formato (\textit{kḷpta}) da un grappolo (\textit{mañjarī}) di fiori di \textit{kadamba}: \textit{Neolamarckia cadamba} (sin. \textit{Anthocephalus cadamba}), albero sacro diffuso in India, caratterizzato da infiorescenze sferiche compatte, profumate, di colore giallo-arancio, simbolo di pienezza e abbondanza;}%
\wordline{ताटङ्क-युगली-भूत-तपन-ओडुप-मण्डला}{tāṭaṅka-yugalī-bhūta-tapana-\\ oḍupa-maṇḍalā}
  { i cui due orecchini (\textit{tāṭaṅka-yugalī}) sono divenuti (\textit{bhūta}) i dischi (\textit{maṇḍala}) del sole (\textit{tapana}) e della luna (\textit{oḍupa}).}%
}

\textit{Śloka} 8 sposta lo sguardo dal volto verso le orecchie, presentando l’elemento ornamentale non come semplice decorazione, ma come forma attraverso cui la luminosità divina diviene articolata e riconoscibile. Il termine \textit{karṇapūra} indica propriamente un ornamento che “riempie” l’orecchio, occupandone lo spazio: non un pendente esterno, ma una presenza che completa e colma l’organo dell’ascolto, suggerendo pienezza e integrazione.

Il secondo \textit{pāda} eleva lo stesso ornamento su un piano cosmico: i due \textit{tāṭaṅka} sono identificati con il sole e la luna. Essi rappresentano energie complementari, testimonianza e misura del tempo, poste ai lati del volto come bilanciamento perfetto. Non si tratta di gioielli semplicemente splendenti, ma di una polarità ordinata che manifesta equilibrio, ritmo e continuità.

Questo \textit{śloka} contiene due \textit{nāma}. \textbf{\textit{Kadamba-mañjarī-kḷpta-\\ karṇapūra-manoharā}} designa l’ornamento auricolare come un grappolo di \textit{kadamba} che colma l’orecchio, sottolineando una bellezza costruita ma vitale. \textbf{\textit{Tāṭaṅka-yugalī-bhūta-tapanōḍupa-maṇḍalā}} presenta la coppia di orecchini come sole e luna, segni di un ordine cosmico reso visibile nella simmetria del volto di Lalitā.
